\documentclass[12pt,a4paper]{article}

\usepackage[utf8]{inputenc}
\usepackage[T1]{fontenc}
\usepackage[margin=2.5cm]{geometry}
\usepackage{hyperref}
\usepackage{microtype}
\usepackage{parskip}
\usepackage{lmodern}

\hypersetup{colorlinks=true, urlcolor=blue!60!black}

\pagestyle{empty}

\begin{document}

\begin{flushright}
Kundai Farai Sachikonye \\
Auf der Lichtung 19, 82178 Puchheim \\
\href{https://github.com/fullscreen-triangle}{github.com/fullscreen-triangle} \\
\texttt{kundai.sachikonye@bitspark.com} \\
(+49) 1626386344 \\[6pt]
16 May 2026
\end{flushright}

\vspace{1em}

Prof.\ Dr.\ Lucie Heinzerling \\
Klinik und Poliklinik für Dermatologie und Allergologie \\
Bereich Dermato-Onkologie \\
LMU Klinikum, Campus Innenstadt, München

\vspace{1.5em}

\textbf{Re: Wissenschaftlicher Mitarbeiter (m/w/d) --- Dermato-Onkologie}

\vspace{1em}

Dear Prof.\ Heinzerling,

I am applying for the scientific staff position in Dermato-Oncology. The core task
you describe --- developing and validating a predictive test for immunotherapy response
--- is the exact problem I have been building a theoretical and computational framework
to address. The framework generates specific, validated predictions about which patients'
tumours will respond to checkpoint inhibitor therapy, derived from first principles
rather than from empirical correlation of clinical variables. I can bring both the
theoretical model and the metabolomics infrastructure to this programme.

\textbf{A predictive test for immunotherapy response.}
My manuscript \textit{On the Thermodynamic Consequences of Categorical Completion
on Biological Systems} derives a structural model of adaptive immune selectivity:
MHC molecules present peptides preferentially from proteins with high categorical
richness $R$ --- quantifying isoform diversity, conformational entropy, and
post-translational modification density. The central empirical result is
$r = 0.73$, $p < 10^{-12}$ across 2,847 IEDB peptides for MHC-I binding
affinity; adding the $R$ term to NetMHCpan raises AUC from 0.68 to 0.81
($p < 0.001$, DeLong test). For melanoma specifically, this generates a
patient-level prediction: tumours with high mutational burden produce neoepitopes
whose categorical richness distribution determines whether anti-PD-1 or
anti-CTLA-4 therapy will generate an effective response --- not just
\textit{whether} the tumour is immunogenic, but \textit{which} checkpoint
is the right target and \textit{which} antigens will expand after treatment.
A companion paper, \textit{On Disease Epistemology in Blind Receiver} (affiliated
with the AIMe Registry and TUM School of Life Sciences Weihenstephan), provides
the binary complement: the loop-holonomy self-consistency criterion achieves
AUC$\,=1.00$ (25/25 validation experiments) as a disease state diagnostic, and
a sparse $\ell_1$ linear programme selects the minimum-side-effect therapeutic
target combination. Together these two results --- a quantitative response
predictor and a binary state classifier --- constitute the architecture of the
predictive test the position describes.
The underlying spectral method is now deployed as four live tools at
\href{https://syndrome-seven.vercel.app}{syndrome-seven.vercel.app}:
an MHC binding predictor (3-channel DFT spectral embedding → 36-dimensional
cosine similarity $\rho$, $O(cL\log L)$, no database lookup);
a neoantigen identifier applying a dual threshold
$\Delta\rho_\text{self} > 0.15$ (immune detection) and $\rho_\text{MHC} > 0.72$
(MHC presentation), with batch ranking by $\rho_\text{MHC} \times \Delta\rho_\text{self}$;
and an immune escape predictor that scans all single-amino-acid substitutions
for variants satisfying $E(v) = -\max\rho(v,\mathrm{Ab}) + w\cdot\rho(v,\mathrm{rec})$
--- antibody evasion with receptor binding maintained.
For melanoma neoantigen vaccine design, you can enter a patient's somatic
mutation directly and receive an immediate prediction of MHC presentation
probability and immune detectability, without a database query and at
$O(cL\log L)$ cost per peptide.

\textbf{Metabolomics.}
My doctoral research at TUM / Bayerisches Forschungsinstitut was in Computational
Lipidomics: mass spectrometry-based analysis of lipid species covering peak
detection, ion annotation against LIPID MAPS and LipidBlast databases, feature
engineering, and multi-omics statistical modelling. Lipidomics is the largest
and most structurally complex sub-discipline of metabolomics, and the analytical
pipelines are identical. At Fraunhofer Institut IGB I worked on production
optimisation of therapeutic lipids; at the Universitätsklinikum Hamburg-Eppendorf
I built automated pipelines for LC-MS/MS and MALDI-TOF characterisation of
proteins and glycans from biological samples.
The \href{https://github.com/fullscreen-triangle/lavoisier}{Lavoisier} framework
I built extends this into a deployable analytical platform: the virtual mass
spectrometry instrument at
\href{https://lavoisier.vercel.app/experiment}{lavoisier.vercel.app/experiment}
runs forward simulation of spectroscopy experiments in the browser, achieving
98.9\,\% feature extraction accuracy against the NIST reference library
and supporting mzML datasets exceeding 100\,GB via distributed computing.
For the microbiome-immunotherapy connection specifically, the relevant metabolite
classes --- short-chain fatty acids, bile acids, indoles, and polyamines --- are
exactly the panel measured in current predictive biomarker studies for checkpoint
inhibitor response in melanoma (Gopalakrishnan et al., Matson et al.), and these
are a standard metabolomics measurement problem. The underlying
\href{https://github.com/fullscreen-triangle/borgia}{Borgia} framework derives
their spectroscopic properties from bounded phase space geometry alone, making it
possible to characterise novel microbial metabolites without calibration curves
against reference standards --- a significant advantage when the relevant
microbiome metabolite panel is not yet fully characterised.

\textbf{Microbiome --- tumour interaction as multi-omics integration.}
The mechanistic link between microbiome composition and immunotherapy outcome
is a data integration problem: metagenomics characterises which microbial species
are present; metabolomics characterises what those species produce;
tumour genomics characterises the antigen landscape; clinical outcomes connect
the three. This is precisely the multi-dimensional data integration problem
the tools I have built are designed for.
The \href{https://github.com/fullscreen-triangle/gospel}{Gospel} framework
processes whole-genome VCF files at $10^6$ variants/second with RNA-seq
integration, covering the tumour genomics and metagenomics sequencing dimensions,
validated against 1000 Genomes Phase~3, gnomAD v3.1.2, and UK Biobank.
Sequence homology search is available directly at
\href{https://vivid-symbolism.vercel.app/search}{vivid-symbolism.vercel.app/search}.
The \href{https://github.com/fullscreen-triangle/syndrome}{Syndrome} framework
integrates these dimensions into a patient-level disease vector
$\mathbf{D} \in [0,1]^8$ --- the genetic expression oscillator ($D_G$) captures
tumour transcriptomic state; the membrane potential oscillator ($D_M$) captures
microbiome-induced barrier integrity changes; the calcium oscillator ($D_{Ca}$)
captures the downstream immune signalling state. The sparse LP for therapeutic
target selection then identifies the minimum-intervention combination ---
which microbiome modulation strategy (probiotics, faecal transplant, antibiotic
preconditioning) combined with which checkpoint inhibitor --- needed to move the
patient's state from the observed disease trajectory to the responder trajectory.

\textbf{Tumour cell spectral diagnostic.}
The \href{https://crown-prince-four-courners.vercel.app}{Crown Prince} platform
implements the Cell Spectral Hologram: given a tumour cell type --- including HeLa
and melanoma lines --- the framework constructs a spectral superposition from
first-principles coupling matrices $K_{ij}$ derived via the Tripartite Isomorphism
(oscillatory dynamics $=$ categorical partition structure $=$ biological membrane
processes). The hologram feeds directly into the loop-holonomy diagnostic, producing
a binary disease state classification and a sparse $\eta^*$ therapeutic perturbation
vector identifying which molecular targets, when modulated, return the cell to the
healthy attractor with minimum intervention. Drug-target binding affinities emerge
as $K_d = \exp(-\Delta M \cdot \ln b)$ from partition geometry, validated against
39/39 NIST binding measurements with zero free parameters --- a 50\,MB representation
covering the same pharmacological space as the 12\,GB DrugBank reference.
This constitutes a second layer of the predictive test architecture: the categorical
richness model (AUC 0.81) identifies which patients will respond to checkpoint
inhibition; the Cell Spectral Hologram (AUC 1.00) identifies which molecular targets
should be perturbed and in what combination.
Live: \href{https://crown-prince-four-courners.vercel.app/tools/spectrum}{crown-prince-four-courners.vercel.app/tools/spectrum}.

\textbf{Background.}
I hold an MSc in Computational Biology (Universität Tübingen) and an MSc in
Pharmaceutical Biotechnology (HAW Hamburg), and completed doctoral research
in Computational Lipidomics at TUM. Published work is affiliated with the AIMe
Registry and TUM School of Life Sciences Weihenstephan. I am legally resident
in Germany with full work authorisation and am available from June 2026.
English is native; German is conversational.

\vspace{1.5em}
Yours sincerely,

\vspace{2.5em}
\textbf{Kundai Farai Sachikonye}

\end{document}
